%% bare_jrnl.tex
%% V1.4b
%% 2015/08/26
%% by Michael Shell
%% see http://www.michaelshell.org/
%% for current contact information.
%%
%% This is a skeleton file demonstrating the use of IEEEtran.cls
%% (requires IEEEtran.cls version 1.8b or later) with an IEEE
%% journal paper.
%%
%% Support sites:
%% http://www.michaelshell.org/tex/ieeetran/
%% http://www.ctan.org/pkg/ieeetran
%% and
%% http://www.ieee.org/
%%
%% ---------------------------------------------------------------
%% Modified by PaperCrew (TFG).
%% Plantilla de referencia: IEEE Journal LaTeX Template oficial
%% (IEEEtran.cls, distribuido vía template-selector.ieee.org),
%% con el front-matter ajustado a las Author Guidelines de
%% IEEE Internet of Things Journal (ieee-iotj.org/guidelines-for-authors).
%% Abstract e introducción se insertan automáticamente.
%% ---------------------------------------------------------------

%%*************************************************************************
%% Legal Notice:
%% This code is offered as-is without any warranty either expressed or
%% implied; without even the implied warranty of MERCHANTABILITY or
%% FITNESS FOR A PARTICULAR PURPOSE!
%% User assumes all risk.
%% In no event shall the IEEE or any contributor to this code be liable for
%% any damages or losses, including, but not limited to, incidental,
%% consequential, or any other damages, resulting from the use or misuse
%% of any information contained here.
%%
%% All comments are the opinions of their respective authors and are not
%% necessarily endorsed by the IEEE.
%%
%% This work is distributed under the LaTeX Project Public License (LPPL)
%% ( http://www.latex-project.org/ ) version 1.3, and may be freely used,
%% distributed and modified. A copy of the LPPL, version 1.3, is included
%% in the base LaTeX documentation of all distributions of LaTeX released
%% 2003/12/01 or later.
%% Retain all contribution notices and credits.
%% ** Modified files should be clearly indicated as such, including  **
%% ** renaming them and changing author support contact information. **
%%*************************************************************************


\documentclass[journal]{IEEEtran}

\hyphenation{op-tical net-works semi-conduc-tor}


\begin{document}

\title{[ARTICLE TITLE]}

\author{[FIRST AUTHOR],~\IEEEmembership{Member,~IEEE,}
        and~[CO-AUTHOR]%
\thanks{[AFFILIATION DATA].}%
\thanks{Manuscript received [DATE].}}

\markboth{IEEE INTERNET OF THINGS JOURNAL,~Vol.~[VOL], No.~[NUM], [MONTH]~[YEAR]}%
{[AUTHORS]: [SHORT TITLE]}

\maketitle

% IEEE IoT-J: el abstract debe tener entre 150 y 250 palabras y no debe
% incluir citas bibliográficas ni ecuaciones.
\begin{abstract}
The Internet of Things (IoT) spans diverse sectors but faces challenges due to technological heterogeneity and costly deployments. This work presents SimulateIoT, a Model-Driven Development (MDD) methodology and toolset for designing, generating, and executing IoT simulations using high-level domain concepts and distributed microservice deployments. The SimulateIoT metamodel formalizes key IoT elements, including Edge, Fog, and Cloud nodes, sensors, actuators, publish-subscribe topics, synthetic and file-based data sources, Complex Event Processing (CEP) engines, NoSQL storage, notifications, and real-time rules. Models created with an EMF/GMF graphical editor and validated by OCL constraints are transformed via Acceleo into Java Thorntail microservices, Mosquitto MQTT brokers, MongoDB databases, CEP engines (Esper, WSO2), and REST APIs. Simulations are orchestrated through Docker Swarm for distributed execution. Two case studies—a smart building and an agricultural environment—demonstrate the platform’s ability to simulate sensor-actuator interactions, CEP rules, data persistence, and system orchestration. The system offers detailed monitoring of CPU, memory, and network resources. SimulateIoT enables scalable, technology-agnostic IoT modeling that streamlines system design, validation, and deployment compared to ad-hoc approaches. Limitations include partial node mobility support, a fixed domain-specific language, simplified hardware abstractions, and static configuration models; planned enhancements aim to introduce dynamic adaptive behaviors. This framework provides a robust foundation for both professional and educational IoT simulation needs.
\end{abstract}

\begin{IEEEkeywords}
[KEYWORD 1], [KEYWORD 2], [KEYWORD 3]
\end{IEEEkeywords}

\IEEEpeerreviewmaketitle


\section{Introduction}









The Internet of Things (IoT) is increasingly applied across diverse domains such as smart cities, agriculture, industry, and home environments. Developing and testing real IoT systems typically demands significant investments in devices, fog and cloud nodes, hardware, and software, as well as considerable time and effort. Consequently, simulation environments have become essential tools to model IoT systems, enabling analysis and optimization without the need for physical deployment. This capability facilitates the development process by reducing costs and complexity.

The state of the art in IoT simulation comprises a broad spectrum of approaches. On the one hand, there are low-level simulators focused primarily on network-level and hardware details. On the other, there are model-driven development (MDE) approaches for IoT systems that enable higher-level abstractions. However, existing solutions generally do not fully integrate complex aspects such as event processing or the architectural layers of Edge, Fog, and Cloud computing at a high abstraction level. This leaves room for approaches that can represent and simulate comprehensive IoT environments with detailed architectural considerations.

Previous works in the area present important limitations that restrict their applicability to comprehensive IoT simulation. Alwasel et al. (2020) [VERIFICAR] propose simulators oriented to big data analysis and cloud IoT environments but lack a graphical interface and have not demonstrated validation in pre-configured environments. Clemente et al. (2018)~\cite{Clemente_2018} apply model-driven development techniques to complex event processing but do not address modeling or simulation of IoT environments. Levis et al. (2003)~\cite{levis2003tossim} developed TinyOS and a simulator focused on low-level mote simulation, lacking a focus on high-level IoT simulation. Mehdi et al. (2014)~\cite{mehdi2014cupcarbon} present simulators for sensor networks emphasizing events and multi-agent models; however, these do not employ model-driven development for system modeling. Zeng et al. (2017)~\cite{Zeng_2017} created IoTSim to analyze IoT applications in cloud contexts but do not provide visual design tools nor support for simulating heterogeneous IoT environments. Papadopoulos et al. (2013)~\cite{papadopoulos2013iotl} enable an experimental platform focused on wireless sensor networks limited to sensors, without full IoT system modeling. Collectively, these approaches cover partial aspects of IoT simulation or model-driven development, yet none offer an integrated visual language and a complete cycle encompassing design, validation, code generation, and deployment of simulated IoT environments with layered Edge, Fog, and Cloud architectures.

This work proposes SimulateIoT, a domain-specific language (DSL) and an accompanying set of tools grounded in model-driven development to address the identified gaps in IoT simulation. The approach enables users to define high-level simulated IoT environments incorporating sensors, actuators, and computing nodes distributed across Edge, Fog, and Cloud layers. By modeling the relationships among components, data flows, complex event processing rules, and storage mechanisms, it facilitates the comprehensive representation of IoT systems.

The use of a graphical editor combined with model-to-text transformations allows automatic generation of deployable artifacts in the form of Docker microservices, which can be instantiated and monitored in real-time. This integrated methodology supports a complete lifecycle—from design and validation to code generation and deployment—thus overcoming limitations of prior work that either lacked visual modeling tools, deployment automation, or did not fully consider the architectural layering inherent in modern IoT systems.

\begin{itemize}
  \item Definition of the SimulateIoT DSL for modeling simulated IoT environments (Sections III and IV).
  \item A metamodel encompassing sensors, actuators, processing nodes, communication topics, and processing rules (Section IV-A).
  \item A graphical editor to visually design IoT environments (Section IV-B).
  \item An automatic code generation engine that produces microservices and supports deployment using Docker Swarm (Section IV-C).
  \item A deployment and monitoring framework for simulated environments incorporating NoSQL databases, MQTT brokers, and Complex Event Processing (CEP) engines (Section IV-D).
  \item Two case studies applied to smart buildings and agricultural environments to validate expressiveness and functionality (Section V).
\end{itemize}

The validation of the proposed approach is conducted through two detailed case studies that model complex IoT environments: a smart university building with multiple sensors, actuators, and Fog/Cloud nodes, and an agricultural environment designed for irrigation using sensors for humidity, pH, temperature, and automated actuators. These case studies demonstrate the capability to generate code and deploy the simulated environments using Docker and microservices, supporting real-time simulation, data processing, monitoring, and the definition of complex event-driven behaviors.

The remainder of the article is organized into several sections that include the introduction, related work, the SimulateIoT methodology, the design and implementation of the tools (covering the metamodel, graphical editor, model-to-text transformations, and deployment), detailed case studies on smart buildings and agricultural environments, followed by a discussion of limitations, concluding remarks, and references.


\section{[SECTION 2 TITLE]}
\label{sec:section2}

[Section 2 content.]


\section{[SECTION 3 TITLE]}
\label{sec:section3}

[Section 3 content.]


\section{[SECTION 4 TITLE]}
\label{sec:section4}

[Section 4 content.]


\section{Conclusion}
\label{sec:conclusion}

[Conclusions.]


\appendices
\section{Appendix}
[Appendix content.]


\section*{Acknowledgment}

[Acknowledgments.]


% Sin bibliografia: no se proporcionaron citas.


\begin{IEEEbiographynophoto}{[FIRST AUTHOR]}
[Brief author biography.]
\end{IEEEbiographynophoto}


\bibliographystyle{IEEEtran}
\bibliography{simulateiot}

\end{document}
